\section{The Whole Action}

\subsection{Reading the sequence as one thing}

Set end to end, the Order of Mass has a shape that no single movement discloses. It opens by
constituting a body: a procession that passes through the people to reach the altar, a corporate
signing, a greeting that requires an answer, an admission of sin made by all, a cry for mercy, a
hymn of praise, and a silence gathered into one petition. It then submits that body to speech it
did not compose: readings it did not choose, a psalm that answers them, a Gospel it stands for,
an exposition, a creed that is not an opinion, and a prayer that turns the body outward toward
the world. Only then does it move to the altar, and it does so by the plainest possible means:
bread, wine, water, a blessing of God for what he has given, a washing of hands, and a request
for prayer. The Eucharistic Prayer is the act; everything before it has been preparation and
everything after it is consequence. Then the body is fed, individually, each person addressed
and each answering. Then it is sent.

Three structural judgements follow from having read the whole in order.

\emph{The rite's centre of gravity is the Eucharistic Prayer, and the rite says so.} The
Instruction calls it the centre and high point of the whole celebration and says of no other
movement anything comparable. Accounts of the reformed Mass that treat the Liturgy of the Word
as its true innovation, or the Communion Rite as its emotional centre, describe a celebration
the book does not legislate.

\emph{The distribution of ministries is the reform's most pervasive structural change, more
pervasive than any textual one.} Lector, psalmist, deacon, cantor, the assembly announcing
intentions, the assembly answering the greeting, saying the \latin{Confiteor}, praying the
\latin{Pater noster}, answering \latin{Amen} at Communion --- almost every movement now has
speakers it did not have. This is what \emph{Sacrosanctum Concilium} 14 asked for and it is the
change a person notices first. It is also the change with the least textual footprint: most of
these parts existed and were spoken by someone else.

\emph{The reform's own self-description is neither ``restoration'' nor ``creation'' but both,
and the distinction is traceable.} Paul VI's constitution names three restorations (homily,
universal prayer, penitential act), one enlargement from tradition and new composition
(prefaces), one addition by decree (three new Canons), one simplification (offertory, fraction,
Communion), and one reorganisation (the Lectionary). This study has followed those categories
because they are the promulgating authority's own and because they can be checked. A reader who
wants to argue that the reform went further than it claimed, or less far, has in that list the
right place to start.

\subsection{Continuity and change, counted}

It is worth stating the balance in the only way that can be checked --- by counting what the
2008 Order of Mass shares with the 1962 Order and what it does not.

Retained with little or no change of wording: the sign of the cross; \latin{Dominus vobiscum}
and its answer; the \latin{Confiteor} in substance; the \latin{Kyrie} in Greek; the
\latin{Gloria}; the Nicene Creed; the mingling formula in substance; \latin{In spiritu
humilitatis}; \latin{Orate, fratres} and \latin{Suscipiat}; the whole preface dialogue; the
\latin{Sanctus}; the Roman Canon from \latin{Te igitur} to \latin{Per ipsum} apart from the
dominical words; the \latin{Pater noster} with its introduction and the embolism; \latin{Pax
Domini}; the commingling formula; the \latin{Agnus Dei}; the priest's two preparation prayers;
\latin{Domine, non sum dignus}; \latin{Quod ore sumpsimus}; \latin{Ite, missa est} and
\latin{Deo gratias}.

This is a historical continuity classification, not a word-level collation. In particular,
the checked target quotation of the prayer for peace ends at \latin{digneris}, so its conclusion
is untested rather than established absent; and the second priestly preparation is titled
\latin{Perceptio Corporis et Sanguinis tui}, whereas the checked older antecedents read
\latin{Perceptio Corporis tui}. No item in this list licenses an older text to stand as the
target where a named divergence or an untested extent remains.

Substantially changed or newly composed: the offertory prayers; the words of consecration (three
specific alterations); the acclamation after the consecration, which is wholly new; the
doxology after the embolism, which is imported; the invitation \latin{Ecce Agnus Dei} in its
present form; three of the four Eucharistic Prayers; and, in 2008, three dismissal formulas.

Changed in manner rather than matter: the Canon is audible; the readings are proclaimed by
others; the people say what the ministers used to say; the Latin remains typical but the
vernacular is normal; and the whole rite is celebrated in a single continuous flow rather than in
parallel streams of said and sung text.

That count is not a verdict. It is the material on which a verdict has to be based, and it is
the reason both of the usual summaries --- ``a new rite'' and ``the same Mass'' --- are too
coarse to be useful.

\subsection{Project synthesis}

\begin{studybox}{Project synthesis --- the judgement this study reaches}
\textbf{The judgement.} The postconciliar Order of Mass, in the state printed in the 2008 Latin
\latin{editio typica tertia emendata}, is best described as a substantially reordered
celebration of a substantially retained rite, in which the deepest textual change is not the
new anaphoras or the reduced offertory but the alteration of the dominical words themselves,
and in which the deepest practical change is not textual at all but the redistribution of
speech among the ministries and the audibility of the Canon. The 2011 English is a faithful
execution of a published translation rule whose internal tension --- between integral rendering
and easy intelligibility --- it inherited rather than created.

\smallskip
\textbf{The reasons.}
\begin{itemize}[itemsep=0.15em]
\item \emph{Retention is demonstrable, not asserted.} The list above of texts carried over from
the 1962 Missal with little or no change covers the great majority of the fixed texts of the
rite, and it includes the Roman Canon in substance. Both books were read directly for this
comparison.
\item \emph{Reordering is equally demonstrable.} The Instruction's own account of the two parts
and their hinge, the transfer of the readings to lectors and deacons, the assembly's assumption
of the \latin{Confiteor}, the \latin{Pater noster}, and the Communion \latin{Amen}, and the
audibility of the Canon are all provisions of the book, not effects of custom.
\item \emph{The consecratory words were changed, and this is the strongest fact on the critical
side.} \latin{Mysterium fidei} was removed from the form of the chalice, \latin{Haec
quotiescumque feceritis} was replaced by \latin{Hoc facite in meam commemorationem}, and
\latin{quod pro vobis tradetur} was added over the bread --- by papal decree, on stated pastoral
and practical grounds, to a form that had stood substantially unchanged in the Latin West for
more than a millennium. Any account of the reform that does not name this is not serious.
\item \emph{The new prayers are not of one kind and cannot be assessed as a bloc.} Prayer II
depends demonstrably, clause by clause, on the Verona Latin anaphora, but supplies from the
Roman tradition a \latin{Sanctus}, a consecratory epiclesis, and intercessions that the ancient
text lacks. Prayer III is a new composition built on Malachi 1:11 with dense sacrificial
vocabulary. Prayer IV is a new composition on the Antiochene plan, structurally alien to Rome
and constrained by its own rubric because of it.
\item \emph{The 2011 English did what its instrument required.} Every principal disputed change
is traceable to \emph{Liturgiam authenticam} 43, 53, or 56, or to the circular letter of 17
October 2006, and the manner of introduction --- all at once, with catechesis --- is what n.\ 74
prescribes.
\end{itemize}

\smallskip
\textbf{The strongest counterargument, at full strength.} The judgement above measures continuity
by counting texts, and counting texts is the wrong measure of a rite. A liturgy is not a corpus
but a form of action, and forms of action can be transformed while every sentence survives. On
this view, the decisive facts are precisely the ones this study classes as ``manner rather than
matter,'' and they are not lesser: a Canon prayed silently by a priest facing east, with the
faithful in interior adherence, is a different act from the same words declaimed audibly across
a table to a congregation, however identical the letters on the page; an offertory whose prayers
name the host and the chalice of salvation forms a different expectation than one that blesses
God for bread; a rite in which the celebrant reads everything himself asserts something about
sacerdotal mediation that a rite of distributed ministries does not. Add the alteration of the
words of consecration --- which this study concedes --- and the enormous permissiveness of the
options (four dismissals, three penitential forms, four Eucharistic Prayers, four chant options
in the United States at three points), and the conclusion follows that the received rite was not
so much retained as quarried: its stones are in the new building, and the building is new. On
this account the study's own \emph{Retained} list is evidence for the counterargument, since a
rite that has to be shown to be continuous by inventory is one whose continuity is no longer
self-evident in the celebrating.

\smallskip
\textbf{Why the judgement is nevertheless retained.} Because the counterargument, at its
strongest, is an argument about celebration and not about the Order of Mass, and it proves less
about the book than it claims. The Missal does not require the postures, orientations, or
performance styles on which the argument's contrasts depend: it does not prescribe celebration
facing the people, it does not forbid singing the Canon, it recommends silence at four points,
and it prints the Roman Canon first. The argument's real target is a widespread manner of
celebrating the reformed rite, and that manner is defensible or not on its own terms. The book's
own centre of gravity, stated in its own Instruction and reflected in its own arrangement, is
where the older book's was.

\smallskip
\textbf{What would change the judgement.} A single finding would do it: documentary evidence,
from the acts of the Consilium or from the promulgating authority, that the changes to the
dominical words were made in order to alter or attenuate what the Church intends by them --- as
distinct from the pastoral and practical grounds Paul VI's constitution states. That would move
the consecratory alteration from the category of textual conformation to the category of
doctrinal revision, and the judgement of substantial retention could not survive it. Short of
that, a demonstration that the reformed Order requires (rather than permits) a manner of
celebration incompatible with the older rite's theology of sacrifice would have the same effect;
this study found no such requirement in the 2008 book, but its reading of the rubrics is not
exhaustive and the General Instruction's chapters III--VIII were not read in full.

\smallskip
\textbf{Status.} This is a source-grounded editorial synthesis by an AI contributor. It is not
magisterial teaching, a canonical opinion, or a specialist consensus, and it has not received
independent liturgical, historical, or theological review.
\end{studybox}

\subsection{What this study does not decide}

It does not decide whether the reform should have been undertaken, whether its execution
honoured \emph{Sacrosanctum Concilium} 23's requirement of organic growth, or what the present
discipline governing the older Missal is. It does not adjudicate among the patristic readings of
\latin{Et cum spiritu tuo}, does not assign authorship to any of the new prayers, and does not
date the sections of the Roman Canon severally. It offers no assessment of the 2011 English of
the Roman Canon, of the orations, or of the Eucharistic Prayers, because assessing them would
require reproducing them. And it makes no judgement about the celebration a reader actually
attends: the distance between the book and the parish is a real and important subject, and it is
a different one.
